%!TEX root = ../main.tex

\begin{abstract}
    Policy-gradient methods are widely used in reinforcement learning, yet training often becomes unstable or slows down as learning progresses. We study this phenomenon through the \emph{noise-to-signal ratio} (NSR) of a policy-gradient estimator, defined as the estimator variance (noise) normalized by the squared norm of the true gradient (signal). Our main result is that, for (i) finite-horizon linear systems with Gaussian policies and linear state-feedback, and (ii) finite-horizon polynomial systems with Gaussian policies and polynomial feedback, the NSR of the REINFORCE estimator can be characterized exactly—either in closed form or via numerical moment-evaluation algorithms—without approximation. For general nonlinear dynamics and expressive policies (including neural policies), we further derive a general upper bound on the variance. These characterizations enable a direct examination of how NSR varies across policy parameters and how it evolves along optimization trajectories (\eg SGD and Adam). Across a range of examples, we find that the NSR landscape is highly non-uniform and typically increases as the policy approaches an optimum; in some regimes it blows up, which can trigger training instability and policy collapse.
\end{abstract}
